\section{Introduction (서론)}

\subsection{연구의 배경: 대형 언어 모델의 시대와 그 이면}
2020년대 이후, Transformer 아키텍처에 기반한 대형 언어 모델(Large Language Models, 이하 LLM)은 기계 학습과 인공지능 역사상 유례없는 파급력을 보여주고 있습니다. GPT-4, Claude 3, Gemini, 그리고 오픈소스인 Qwen 및 Llama 시리즈에 이르기까지, 최신 모델들은 방대한 양의 텍스트 데이터를 섭취하여 인간의 언어를 능숙하게 구사하게 되었습니다. 이들은 문서를 요약하고, 창의적인 시를 작성하며, 프로그래밍 언어로 어플리케이션을 코딩하는 등 지식 노동의 상당 부분을 자동화하는 수준에 도달했습니다.

그러나 이러한 언어적 유창함의 이면에는 치명적인 약점이 존재합니다. 바로 `엄밀한 논리와 수학적 계산(Strict Logic and Mathematical Calculation)' 능력의 부재입니다. 일반적인 자연어 처리 작업과 달리 수학 문제는 단 한 번의 연산 실수나 논리적 비약이 최종 결과의 완전한 붕괴를 초래합니다. "12345와 67890을 곱하면 얼마인가?" 혹은 "$100!$ 의 끝에 연속된 $0$의 개수는 몇 개인가?"와 같은 질문에 대해, LLM은 정답이 아닌 '그럴듯한 오답'을 매우 자신감 있는 어조로 출력하는 현상, 즉 환각(Hallucination)을 빈번하게 발생시킵니다.

\subsection{AI는 왜 수학을 못하는가? (The Mathematical Bottleneck)}
LLM이 수학 문제에 취약한 원인은 모델의 근본적인 작동 방식에서 기인합니다. 

첫째, \textbf{자기 회귀(Auto-regressive) 특성으로 인한 에러 누적}입니다. LLM은 이전까지 생성된 토큰(단어의 조각)들을 바탕으로 바로 다음 토큰의 확률을 계산하여 텍스트를 이어나갑니다. 만약 추론의 첫 번째 단계에서 사소한 덧셈 실수를 저지르게 되면, 모델은 그 오답을 기정사실화하여 이후의 모든 논리를 그 오답에 맞춰 전개하게 됩니다. 이를 확증 편향(Confirmation Bias)이라 하며, 스스로 실수를 인지하고 되돌아가는 능력(Backtracking)이 부족하기 때문입니다.

둘째, \textbf{토크나이저(Tokenizer)의 한계}입니다. 언어 모델은 텍스트를 의미 단위로 쪼개어 인식하는데, 숫자의 경우 그 패턴이 불규칙하게 쪼개지곤 합니다. 예를 들어 '12345'라는 숫자가 '123'과 '45'로 분리되어 벡터 공간에 저장되면, 숫자의 자릿수와 물리적 크기에 대한 감각(Number Sense)이 상실됩니다. 이로 인해 모델은 숫자의 덧셈과 곱셈을 진정한 의미의 산술(Arithmetic)이 아닌 텍스트 패턴 매칭(Pattern Matching)으로 해결하려 시도하게 됩니다.

셋째, \textbf{작업 기억(Working Memory)의 부재}입니다. 인간은 복잡한 다단계 연산을 수행할 때 종이에 중간 결과를 적어두거나(Scratchpad), 외부의 계산기를 활용합니다. 반면 기본 상태의 LLM은 모든 연산을 오로지 신경망 내부의 파라미터 연산(머릿속 암산)만으로 해결해야 하는 가혹한 제약 조건 하에 놓여 있습니다.

\subsection{OMI (Orchestrated Math Interpreter) 제안: 발상의 전환}
본 연구는 "AI에게 머릿속으로만 계산하도록 강요하는 것이 옳은가?"라는 질문에서 출발했습니다. 인간 수학자가 어려운 문제를 풀 때 계산기를 사용하고 공학용 소프트웨어를 활용하듯, AI에게도 동일한 도구를 쥐여준다면 어떨까요?

이러한 철학을 바탕으로 개발된 \textbf{OMI (Orchestrated Math Interpreter)} 시스템은 AI의 역할을 단순히 '답을 뱉어내는 자'에서 '문제를 어떻게 풀지 계획하고, 필요한 도구를 사용해 답을 조립하는 지휘자(Orchestrator)'로 재정의합니다. OMI 시스템 내에서 LLM은 파이썬(Python)이라는 강력하고 결함 없는 산술 연산 엔진과 결합됩니다. AI는 자연어로 문제를 분석하여 해결 논리를 파이썬 코드로 작성하고, 샌드박스(Sandbox) 환경에서 코드를 실행하여 나온 정확한 결과값을 토대로 수학적 결론을 도출합니다. 

이를 통해 AI의 가장 큰 약점이었던 '연산 실수'를 파이썬 인터프리터가 원천 차단하고, AI는 본연의 강점인 '논리적 방향성 제시'와 '코드 작성'에만 집중할 수 있게 됩니다.

\subsection{주요 기여도 및 연구 성과 (Key Contributions)}
본 프로젝트는 단순히 파이썬을 붙이는 수준을 넘어, 수학 교육 및 풀이 시스템의 패러다임을 바꿀 수 있는 혁신적인 파이프라인과 데이터셋 생성 기법을 포함하고 있습니다. 주요 기여도는 다음과 같습니다.

\begin{enumerate}
    \item \textbf{5단계 심층 추론 파이프라인(5-Stage Deep Reasoning Pipeline) 구축:} 
    문제를 도메인별로 라벨링(Labeling)하고, 맞춤형 지식을 검색(Retrieval)하며, 문제 스케일에 따라 시뮬레이션 혹은 이론적 접근으로 경로를 나누는 라우팅(Routing), 그리고 안전한 파이썬 실행(Execution)과 축소 검증(Verification)으로 이어지는 총체적 워크플로우를 확립했습니다.
    
    \item \textbf{다중 에이전트 앙상블(Multi-Agent Ensemble) 기반의 무결성 확보:} 
    패스 앳 케이(Pass@k) 확률 모델을 활용한 다수결 투표(Majority Voting) 시스템과, Coder와 Reviewer 역할을 분리하여 AI의 확증 편향을 끊어내는 다중 에이전트 토론(Debate) 메커니즘을 도입하여 난제에 대한 정답률을 극대화했습니다.
    
    \item \textbf{역공학(Reverse Engineering) 기반의 무오류 합성 데이터 엔진:} 
    데이터 오염(Data Contamination)을 막고 완벽한 참(True)이 보장되는 고품질 훈련 데이터를 무한대로 생성하기 위해, 정답으로부터 수식을 역으로 파생시키고 이를 자연어 문제로 감추는(Obfuscation) 새로운 데이터 합성 기법(Open-Reverse-Math)을 제안합니다.
    
    \item \textbf{Glass-box 기반의 교육용 웹 대시보드 및 XAI 시각화:} 
    백엔드에서 벌어지는 복잡한 라우팅과 파이썬 실행 과정을 프론트엔드 대시보드(Vite+React)를 통해 실시간으로 렌더링합니다. 또한 XAI(Explainable AI) 기술을 통해 AI의 전략 선택 이유를 인간의 언어로 풀어내어 투명하고 신뢰할 수 있는 교육 시스템을 완성했습니다.
\end{enumerate}

\subsection{보고서의 구성 (Document Structure)}
본 보고서는 총 5개의 섹션으로 구성되어 있으며, 각각의 섹션은 OMI 시스템의 핵심 요소들을 심도 있게 다룹니다.
제 2장(Related Work)에서는 Tool-Integrated Reasoning과 다중 에이전트 시스템에 관한 학계의 기존 연구와 그 한계를 분석합니다.
제 3장(Methodology)에서는 본 연구의 핵심인 5단계 심층 추론 파이프라인의 알고리즘 원리와 수학적 증명, 그리고 역공학 데이터 생성 엔진의 상세 로직을 묘사합니다.
제 4장(System Implementation)에서는 Vertex AI 기반의 클라우드 아키텍처와 이를 시각화하는 대시보드의 실전 구동 시나리오를 다룹니다.
마지막으로 제 5장(Conclusion)에서는 본 프로젝트의 의의와 향후 확장성에 대해 논의하며, 부록(Appendix)을 통해 비전공자도 쉽게 이해할 수 있는 상세한 용어 사전을 제공합니다.
